\section{Explainability Analysis}\label{sec:xai}

A high test accuracy is necessary but not sufficient evidence that a
medical-imaging classifier is clinically defensible: the model must also be
shown to attend to lesion tissue rather than to incidental cues such as teeth,
lip outlines, or lighting artifacts. To supply that evidence, every model in
the study is accompanied by two complementary explanations for both heads:
a Gradient-weighted Class Activation Mapping++ (Grad-CAM++) saliency panel and
a Local Interpretable Model-Agnostic Explanations (LIME) boundary-mask panel.
The artifacts are produced by a single script, \texttt{explain\_model.py}, with
identical preprocessing across all ten architectures, and are stored per model
as \texttt{explain\_binary.png} and \texttt{explain\_subtype.png}. This section
reports the patterns observed on the proposed model
(Section~\ref{subsec:gradcampp} and Section~\ref{subsec:lime}) and then ranks
all models by how closely their saliency aligns with visible pathology
(Section~\ref{subsec:xai_compare}).

\subsection{Grad-CAM++ Saliency}\label{subsec:gradcampp}

Grad-CAM++~\cite{chattopadhyay2018gradcampp} extends
Grad-CAM~\cite{selvaraju2017gradcam} by improving the localization of multiple,
and of spatially small, instances of the target class: it replaces the global
gradient average with a pixel-wise weighted combination of higher-order
gradients. We adopt it as the gradient-based saliency method for every backbone,
computed on the final stage of the network so that the map reflects the
representation immediately consumed by the classification heads:

\begin{itemize}
  \item For the convolutional baselines (ResNet50, DenseNet121,
  ConvNeXt-Tiny, the EfficientNet variants, and Inception~V3) the target layer
  is the last convolutional block.
  \item For Swin-T the target is the last attention block prior to global
  pooling.
  \item For the proposed Custom EfficientNet~V2 the target is Stage~5, the
  MBConv+SE block immediately downstream of the AttentionHub, so that the
  map captures the feature representation the hub hands to Stage~5.
\end{itemize}

\noindent For each image the class index is taken from the relevant head's
predicted argmax, and the resulting per-image attention map is overlaid on the
input.

Inspection of the proposed model's panels (\texttt{explain\_binary.png} and
\texttt{explain\_subtype.png} of the \texttt{custom\_efficientnet\_v2\_hub\_v2}
run) reveals four recurring patterns. \emph{First}, lip-margin benign cases
receive broad, diffuse attention spread across the visible mucosal area with no
focal fixation on the teeth, indicating that the model reasons over the
lesion-bearing tissue rather than over incidental hard-tissue cues.
\emph{Second}, white-plaque benign cases receive a single tight focal spot
placed directly on the visible plaque; this is the most spatially concentrated
activation in the panel and is anatomically correct. \emph{Third},
diffuse-mucosa malignant cases receive a distributed, multi-spot activation
over the lesion region, which is appropriate for lesions whose boundary is not
sharply defined. \emph{Fourth}, focal-bump malignant cases receive two crisp
focal spots that coincide with the two visible bumps, in clinical terms the
kind of annotation a junior clinician would draw on the image directly.

A clarification on what this map does and does not probe is warranted. The
AttentionHub is located at Stage~4, whereas the saliency above is computed at
Stage~5, one block downstream of the hub. The map therefore characterises the
representation that the classification heads consume, not the internal
behaviour of the hub itself; we make no causal attribution of these Stage-5
patterns to specific hub modules. What the map does support is the property
that matters clinically: the Stage-5 representation handed to the heads is
lesion-focused, exhibiting \emph{broad} attention on diffuse lesions and
\emph{focal} attention on bump-like lesions and showing no fixation on teeth or
lip outlines. This is qualitatively consistent with the role-complementarity
argument that motivated the v2 design in Section~\ref{sec:ablation}, without
being treated as direct evidence of it. Figure~\ref{fig:xai_proposed} collects
six hand-curated test cases that exhibit these patterns.

\begin{figure*}[!t]
  \centering
  \includegraphics[width=\textwidth]{fig06c_proposed_explain_panel_landscape}
  \caption{Explainability showcase for the proposed Custom EfficientNet~V2
  (AttentionHub-v2). Six curated test cases are arranged in a $3\times2$ grid;
  within each cell the input is shown next to its Grad-CAM++ heat-map and its
  LIME boundary mask (columns \emph{Original} $\mid$ \emph{Grad-CAM++} $\mid$
  \emph{LIME}). The cases span the binary head (a malignant ulcerated lesion)
  and the subtype head (CoS lip and perioral-skin lesions, an OC raised tongue
  lesion, a Gum inflamed lip margin, and OLP striations). Across all cases the
  Grad-CAM++ peak coincides with the visible lesion, showing a tight focal spot on
  compact lesions and an elongated map tracing the contour of linear lesions,
  and the LIME contour traces the lesion boundary rather than surrounding
  teeth or skin.}
  \label{fig:xai_proposed}
\end{figure*}

\subsection{LIME Boundary Masks}\label{subsec:lime}

Grad-CAM++ is gradient-based and therefore tied to the internal structure of
the network; to obtain a model-agnostic second opinion we additionally compute
LIME~\cite{ribeiro2016lime}. LIME segments each test image into super-pixels,
perturbs them by occlusion, and fits a sparse interpretable surrogate to the
resulting changes in predicted class probability, yielding a boundary mask over
the super-pixels whose perturbation most strongly affects the decision. Because
LIME treats the classifier as a black box, agreement between its masks and the
Grad-CAM++ heat-maps is non-trivial corroboration rather than a restatement of
the same signal.

For the proposed model the two explanations are mutually consistent: the LIME
masks align with the Grad-CAM++ heat-maps and with the visible lesion tissue
across the curated cases of Figure~\ref{fig:xai_proposed}. The agreement is
sharpest on the focal-bump cases, where the LIME boundaries are tight around
the bumps themselves. This property is not universal across the benchmark:
several baselines, notably ConvNeXt-Tiny under from-scratch training and Swin-T
on small lesions, produce LIME boundaries that bleed into the teeth or the
surrounding skin. The convergence of a gradient-based and a perturbation-based
explainer on the same lesion-centred regions strengthens the claim that the
proposed model's decisions rest on pathology rather than on surrounding context.

\subsection{Cross-Model Comparison}\label{subsec:xai_compare}

To place the proposed model's interpretability in context, the Grad-CAM++ maps
of the strongest architectures are compared on a fixed set of input images.
Because all runs share the same deterministic seed and split, a given row index
in \texttt{explain\_binary.png} corresponds to the same physical image across
models, so the comparison isolates the effect of the architecture rather than
of the input. Figure~\ref{fig:xai_cross} arranges these maps for the proposed
model and four representative baselines on four common binary samples.

A side-by-side review of the Grad-CAM++ panels supports a qualitative
three-tier ordering, from most to least clinically aligned. We emphasise that
this ordering is a visual judgement of the released panels rather than a
quantitative interpretability measurement: the dataset carries no lesion masks,
so mask-based metrics such as pointing-game accuracy cannot be computed, and no
quantitative interpretability score is reported here. The ordering should
therefore be read as indicative. The \emph{most aligned} tier comprises the
Custom EfficientNet~V2 (AttentionHub-v2), EfficientNet~V2-B2, and
Inception~V3: all three concentrate their activation on lesion tissue with
negligible response on teeth or skin. The \emph{mid-tier} comprises ResNet50,
DenseNet121, and EfficientNet-B0, which generally localize the lesion correctly
but exhibit occasional spurious activation on lip outlines. The \emph{weakest}
tier is ConvNeXt-Tiny, whose attention often broadens across the entire oral
cavity without lesion-specific concentration, a behaviour consistent with
its low quantitative accuracy (95.32\,\% binary, 90.46\,\% subtype, see
Table~\ref{tab:benchmark}) and with its failure to reach the early-stopping
criterion under from-scratch training.

The interpretability advantage of the proposed model is therefore not merely a
quantitative claim derived from accuracy: it is directly visible in the
Grad-CAM++ and LIME panels released with this paper. The same architecture that
attains the highest subtype accuracy in the benchmark among the nine baselines
also produces the saliency maps most tightly bound to visible pathology, and a
gradient-based and a model-agnostic explainer agree on this lesion-centred
alignment. This evidence is qualitative; a quantitative interpretability
measure is identified as future work in Section~\ref{subsec:limitations}.

\begin{figure*}[!t]
  \centering
  \includegraphics[width=\textwidth]{fig06_gradcam_cross_model_composite}
  \caption{Cross-model Grad-CAM++ comparison on four common binary test samples
  (two benign, two malignant). Rows correspond to ResNet50, DenseNet121,
  EfficientNet~V2-B2, Inception~V3, and the proposed Custom EfficientNet~V2
  (AttentionHub-v2, highlighted); columns correspond to the four shared inputs.
  The deterministic seed and split guarantee that each column is the same
  physical image across all rows, so differences in the heat-maps reflect the
  backbone alone. The proposed model and the two strongest baselines
  (EfficientNet~V2-B2, Inception~V3) concentrate activation on lesion tissue,
  whereas weaker backbones spread activation onto lip outlines and surrounding
  regions.}
  \label{fig:xai_cross}
\end{figure*}
